%Prépa Agreg

\setcounter{equation}{0}
\setcounter{Q}{0}

Le diagramme ci-dessous correspond au binaire \ce{Au-Si}. On considèrera une solution liquide idéale et les $\Delta_{fus} H^\circ$ constants au voisinage de l'eutectique.

\begin{center}
\includegraphics[width=0.5\textwidth]{figures/Fig01}
\captionof{figure}{Diagramme isobare \ce{Au-Si} à une pression de \SI{1}{\bar}}
\end{center}

\Q Déterminer $\Delta_{fus} H^\circ (\ce{Au})$ et $\Delta_{fus} H^\circ (\ce{Si})$ au voisinnage de la température de l'eutectique $T_E$.

\Q  Comparer aux $\Delta_{fus} H^\circ$ à la température de fusion : 
\begin{align*}
\Delta_{fus} H^\circ (\ce{Au}) &= 12.68 ~ \si{\kilo\joule\per\mole} \textbf{ à} ~ 1064.4 ~ \si{\celsius} \\
\Delta_{fus} H^\circ (\ce{Si}) &= 39.62  ~ \si{\kilo\joule\per\mole} \textbf{ à} ~ 1410 ~ \si{\celsius}
\end{align*}
Commenter.


\paragraph*{Données} $R = 8.314 ~ \si{\joule\per\mole\per\kelvin}$
On rappelle que :
\begin{align*}
\frac{\partial \frac{\mu}{T}}{\partial T} = - \frac{H_m}{T^2}
\end{align*}

avec $\mu$ le potentiel chimique d'une phase et $H_m$ l'enthalpie molaire d'une phase.